\section{Экспериментальная часть}

Постановка задачи, обозначения и исследовательские вопросы --- в главе 1 (§1.2, §1.6--1.8); теоретическая база --- в главе 2.

\subsection{Цель и постановка}\label{ux446ux435ux43bux44c-ux438-ux43fux43eux441ux442ux430ux43dux43eux432ux43aux430}

Цель экспериментальной части --- эмпирически исследовать онлайн-эксперимент в рекомендательной постановке (выбор варианта показа, метрика CTR): сначала проверить, приносит ли адаптивное назначение трафика измеримую пользу по сравнению с фиксированным разбиением трафика, затем --- какие процедуры дают корректную интерпретацию логов и статистически обоснованный вывод о превосходстве варианта B. Постановка согласуется с OBD \cite{ref26} и типичными задачами маркетплейсов; в других доменах целевой метрикой может быть конверсия, выручка, удержание и т. д. Проверяются RQ1--RQ7 (§1.6).

В индустрии эти вопросы часто смешивают: высокий CTR во время MAB-теста подают как доказательство превосходства B над A. Глава 3 разводит их явно. Блок I (польза онлайн-эксперимента) отвечает: даёт ли адаптивное назначение больше кликов и меньший regret, чем \texttt{fixed\_ab}? Блок II (статистическая значимость и интерпретация логов) отвечает: можно ли по логам того же эксперимента обоснованно заявить «B лучше A» --- и если нет, какие известные процедуры это исправляют?

В рамках выбранной метрики удобно разделить два родственных, но разных вопроса. Первый --- статистический: после эксперимента можно ли подтвердить гипотезу о превосходстве варианта B над A, если нас интересует рост CTR у тестового варианта по сравнению с контрольным. Для такого вывода обычно запускают A/B-тест с фиксированным разбиением трафика, чаще всего 50/50. Второй --- операционный: как не терять ценные события (в нашей постановке --- клики), пока эксперимент ещё идёт, не отдавая половину показов заведомо слабому варианту. Здесь применяют бандитов, которые перенаправляют трафик к варианту с лучшим промежуточным CTR. Первый вопрос относится к статистическому выводу: корректность p-value, контроль ошибки I рода. Второй --- к онлайн-оптимизации: сколько кликов накоплено за горизонт. Ответ на один не подменяет ответ на другой.

\textbf{Структура главы по двум блокам:}

\begingroup\footnotesize\setlength{\tabcolsep}{4pt}
\begin{xltabular}{\linewidth}{@{}
  >{\raggedright\arraybackslash}Y
  >{\raggedright\arraybackslash}Y
  >{\raggedright\arraybackslash}Y
  >{\raggedright\arraybackslash}Y@{}}
\toprule\noalign{}
\midrule\noalign{}
\endhead
\bottomrule\noalign{}
\endlastfoot
\textbf{I. Польза адаптивного онлайн-эксперимента} & MAB / контекстный бандит снижают regret и дают больше кликов? & E1, E6, E8b, E12; абляции & онлайн-оптимизация \\
\textbf{II. Статистическая значимость и логи} & Как корректно интерпретировать логи и получить вывод о превосходстве варианта B? & E4, E5, E11, E14; статистический вывод в E12 & статистический вывод \\
\emph{Дополнение} & Как интерпретировать уже собранные логи без нового трафика? & E2, E2b/c, E9, E13 & OPE \\
\emph{Инженерия} & Успевает ли алгоритм в продакшене по времени задержки? & P1 & применимость \\
\end{xltabular}
\endgroup

Блок I не доказывает статистическую значимость; блок II не измеряет regret. OPE и P1 дают контекст для полной картины онлайн-эксперимента, но центральная линия главы --- связка «польза → проблема статистического вывода → известные решения».

Обозначения политик и метрик --- в §1.2.

Эксперименты реализованы на Python в пакете \texttt{src/}; корректность проверяется автоматическими тестами. Политика \texttt{fixed\_ab} моделирует только схему распределения трафика, а не полный цикл A/B-теста с расчётом размера выборки --- последний вынесен в блок II (E4, E5, E11, E14).

\subsection{Четыре типа оценки}\label{ux447ux435ux442ux44bux440ux435-ux442ux438ux43fux430-ux43eux446ux435ux43dux43aux438}

Полное описание четырёх типов оценки (онлайн-оптимизация, статистический вывод, OPE, продуктовая применимость) --- в §1.8. В главе 3 эксперименты сгруппированы по этим типам; центральная линия --- блоки I и II (§3.1). Смешивать метрики оптимизации и статистического вывода в одном заключении методологически некорректно \cite{ref5,ref21}.

\subsection{Исходные данные}\label{ux438ux441ux445ux43eux434ux43dux44bux435-ux434ux430ux43dux43dux44bux435}

\textbf{Синтетическая среда} моделирует K вариантов с известными вероятностями клика. В E1 --- пять рук с CTR от 3\% до 5\%; в E12 --- две руки с CTR 5\% и 8\%. На каждом шаге политика выбирает вариант и получает бинарный отклик. Синтетика позволяет однозначно интерпретировать regret и проводить абляции (E6 по параметру $\varepsilon$).

\textbf{Open Bandit Dataset (OBD)} --- логи fashion e-commerce ZOZOTOWN \cite{ref26}. Каждая запись --- показ товара: \texttt{item\_id}, позиция, признаки пользователя, клик и propensity. Подвыборки обозначаются как \texttt{политика\_логирования/сегмент}: первая часть --- как собирали лог (\texttt{random} --- равномерное случайное назначение товаров; \texttt{bts} --- адаптивная политика BTS, адаптивное логирование); вторая --- аудитория (\texttt{all} --- все пользователи; \texttt{men}, \texttt{women} --- по полу). Например, \texttt{random/all} --- лог при случайном показе по всей аудитории; \texttt{bts/all} --- лог при адаптивном назначении по всей аудитории. В курсовой используются учебная подвыборка small OBD (10 000 событий из \texttt{random/all}, 38--46 кликов, CTR около 0,4\%) и полный \texttt{random/all} (1 374 327 событий, 4 768 кликов, E9a). В кампании \texttt{all} --- 80 товаров; это многорукая постановка, а не классическая A/B с двумя вариантами.

\subsubsection{Разведочный анализ OBD (EDA)}\label{ux440ux430ux437ux432ux435ux434ux43eux447ux43dux44bux439-ux430ux43dux430ux43bux438ux437-obd-eda}

Перед OPE и проекцией на пару вариантов (E13) выполнен разведочный анализ учебной подвыборки \texttt{random/all} (10 000 событий). По полю \texttt{item\_id} агрегированы показы, клики и CTR; в кампании all зафиксировано 80 уникальных товаров. Суммарный CTR подвыборки --- 0,38--0,46\% (38--46 кликов на 10 000 показов): типичный для fashion e-commerce низкий CTR, при котором даже миллион событий даёт лишь тысячи кликов (E9a).

Распределение CTR по товарам сильно скошено: у большинства позиций в топ-20 по показам CTR от 0\% до 1,5\%; единичные клики на сотни показов --- норма, а не выброс. На рис. 3.5 показаны CTR двадцати самых частых \texttt{item\_id}; разброс подтверждает, что ранжирование алгоритмов по SNIPS на таком объёме будет шумным (E2, E9).

Для попарной постановки E13 (две руки A и B) из топ-20 по показам (\texttt{min\_impressions\ =\ 100}) выбрана пара \texttt{item\_id} 41 и 50 с минимальным \textbar CTR\_A - CTR\_B\textbar{} --- квази-null: по 136 показов, CTR 0,735\% у обоих, gap 0. Артефакты: \texttt{outputs/obd\_pair/item\_ctr\_stats.csv}, \texttt{pair\_selection.csv}.

\begin{figure}[htbp]
\centering
\includegraphics[width=0.92\textwidth]{fig_3_5_obd_ctr_top20.png}
\caption{Рисунок 3.5 — OBD: CTR топ-20 товаров по числу показов (10k \texttt{random/all})}
\label{fig:рисунок_3_5___obd__ctr_топ_20_товаров_по}
\end{figure}

\textbf{Классическая A/B-постановка} (E12, E13) моделирует стандартный дизайн: контрольный вариант A и тестовый вариант B, гипотеза CTR\_B \textgreater{} CTR\_A. E12 решает это на чистой синтетике; E13 проецирует логи OBD на две выбранные \texttt{item\_id} (§3.6.4).

\textbf{Ограничения:} один домен, низкий CTR, упрощение интерфейса (одна позиция). Связка «синтетика + OBD» даёт контролируемый бенчмарк и реальные логи с propensity.

\subsection{Метрики}\label{ux43cux435ux442ux440ux438ux43aux438}

В главе используются четыре типа оценки (§1.8). \textbf{Онлайн-оптимизация} --- измерение пользы адаптивного эксперимента \emph{в ходе} его проведения (не «онлайн» в смысле интернета): суммарное число кликов за горизонт (в таблицах --- \texttt{reward} или \texttt{cumulative\_reward}), regret и \texttt{suboptimal\_share} --- доля шагов, на которых показан не лучший вариант (см. §1.2, сокращения).

\textbf{Статистический вывод} --- проверка гипотезы о превосходстве B над A: rejection rate (доля прогонов с p-value \textless{} $\alpha$), p-value, доверительный интервал, бинарное решение об отклонении H$_0$. \textbf{Последовательный групповой дизайн} (OBF, см. сокращения) и корректный мониторинг при непрерывном наблюдении (mSPRT) --- способы контролировать ошибку I рода при подглядывании за промежуточными результатами фиксированного A/B; это отдельно от адаптивного назначения бандита (§3.6.2).

При OPE --- SNIPS-оценка CTR политики-кандидата, доля совпадений с логом и ESS.

При оценке применимости --- среднее время одного решения алгоритма в микросекундах.

\subsection{Программная реализация}\label{ux43fux440ux43eux433ux440ux430ux43cux43cux43dux430ux44f-ux440ux435ux430ux43bux438ux437ux430ux446ux438ux44f}

Политики --- в \texttt{bandits/}; среды --- в \texttt{environments/}; IPS/SNIPS replay-OPE --- в \texttt{ope/}; процедуры A/B-вывода --- в \texttt{ab\_testing/} (\texttt{inference.py}, \texttt{sequential.py}, \texttt{bandit\_logs.py}, \texttt{weighted\_inference.py}). Ниже результаты сгруппированы по четырём типам оценки (§1.8, §3.2); внутри каждого блока эксперименты идут по кодам E1--E15, затем P1. Соответствие RQ и экспериментов --- в табл. 1.2 (§1.6); сводка параметров прогонов --- в табл. 3.1.

\textbf{Таблица 3.1 --- Сводка экспериментов}

В столбце «Тип оценки» --- один из четырёх типов по §1.8: онлайн-оптимизация (польза в ходе эксперимента), статистический вывод (проверка гипотезы), OPE (оценка по логам), продуктовая применимость.

\begingroup\small\setlength{\tabcolsep}{5pt}\renewcommand{\arraystretch}{1.12}
\begin{tabularx}{\linewidth}{@{}>{\raggedright\arraybackslash}p{1.05cm}>{\centering\arraybackslash}p{0.95cm}>{\raggedright\arraybackslash}>{\hsize=1.2\hsize}X>{\raggedright\arraybackslash}>{\hsize=1.0\hsize}X>{\raggedright\arraybackslash}>{\hsize=0.8\hsize}X@{}}
\toprule
Код & RQ & Тип оценки & Данные & Объём \\
\midrule
E1 & RQ1 & онлайн-оптимизация & синтетика, 5 рук & T=5000, 20 seed \\
E2 & RQ2 & OPE replay (IPS/SNIPS) & OBD \texttt{random/all} & 10 000 соб., 20 seed \\
E2b/c & RQ6 & OPE replay (IPS/SNIPS) & OBD ×4 подвыборки & 10 000 соб., 20 seed \\
E4 & RQ3.1 & статистический вывод, Monte Carlo & синтетика, 2 руки & 200 trials, T=20 000 \\
E5 & RQ3.1 & статистический вывод, иллюстрация & синтетика, 2 руки & T=5000, 1 seed \\
E6 & RQ4 & онлайн-оптимизация, абляция $\varepsilon$ & синтетика, 5 рук & T=5000, 10 seed ×3 $\varepsilon$ \\
E8b & RQ5 & онлайн-оптимизация, контекстный бандит & синтетика & T=10 000, 20 seed \\
E9a & RQ2 & OPE, потоковая обработка & OBD полный \texttt{random/all} & 1,37 млн соб., 10 seed \\
E9b & RQ6 & OPE, потоковая обработка & OBD \texttt{bts/all} & 1 млн соб., 10 seed \\
E11 & RQ3.2 & статистический вывод, IPS & синтетика, 2 руки & 100 trials, T=20 000 \\
E14 & RQ3.3 & последовательный вывод / подглядывание & синтетика, 2 руки & 200 trials, T=20 000 \\
E15 & --- & статистический вывод, CUPED & синтетика, 2 руки & 100 trials, T=20 000 \\
E12 & RQ7 & онлайн-оптимизация + статистический вывод & синтетика, 2 руки & T=10 000, 20 seed \\
E13 & RQ7 & OPE (иллюстрация) & OBD, проекция на пару & 272 соб., приложение А \\
P1 & продукт & время задержки & --- & 20 000 вызовов, 80 рук \\
E3, E8 & --- & разведочный пакетный прогон (приложение) & OBD & см. приложение А \\
\bottomrule
\end{tabularx}
\endgroup

\subsection{Эксперименты и результаты}\label{ux44dux43aux441ux43fux435ux440ux438ux43cux435ux43dux442ux44b-ux438-ux440ux435ux437ux443ux43bux44cux442ux430ux442ux44b}

Эксперименты сгруппированы по двум центральным блокам (§3.1). §3.6.1 и начало §3.6.4 (онлайн-часть E12) относятся к блоку I: измеряем пользу адаптивного онлайн-эксперимента. §3.6.2 и часть E12 про статистический вывод --- к блоку II: проверяем, как получить статистически корректный вывод и где дают сбой наивные процедуры. §3.6.3 (OPE) и §3.6.5 (P1) дополняют картину интерпретации логов и инженерной применимости.

\subsubsection{Блок I: польза адаптивного онлайн-эксперимента (E1, E6, E8b)}\label{ux431ux43bux43eux43a-i-ux43fux43eux43bux44cux437ux430-ux430ux434ux430ux43fux442ux438ux432ux43dux43eux433ux43e-ux43eux43dux43bux430ux439ux43d-ux44dux43aux441ux43fux435ux440ux438ux43cux435ux43dux442ux430-e1-e6-e8b}

Классический A/B-тест фиксирует разбиение трафика на весь срок: это корректно для проверки гипотезы, но расточительно для метрик --- часть трафика всё время уходит на слабый вариант. Бандит максимизирует CTR во время эксперимента, смещая трафик к лучшему варианту. Ниже проверяется, насколько выигрыш заметен на контролируемой синтетике.

\textbf{Эксперимент E1 (RQ1).} Пятирукая Bernoulli-среда с CTR от 3\% до 5\%; оптимальный вариант --- 5\%. Горизонт T = 5000, 20 seed. Сравниваются \texttt{fixed\_ab} (равномерно по 20\% на руку), $\varepsilon$-greedy ($\varepsilon$ = 0,1), UCB1 и Thompson Sampling. Для \texttt{fixed\_ab} при K = 5 теоретически suboptimal\_share = 1 - 1/K = 0,80: одна рука оптимальна, четыре --- нет. В табл. 3.2 эмпирическое значение 0,80 совпадает с расчётом. У адаптивных политик suboptimal\_share ниже (0,63--0,77), regret --- меньше.

\textbf{Таблица 3.2 --- E1: онлайн-оптимизация (среднее ± ст. откл. по 20 seed, T = 5000)}

Столбцы: \texttt{policy\_name} --- политика; \texttt{reward} --- суммарное число кликов; \texttt{regret} --- недобранные клики относительно оракула; \texttt{suboptimal\_share} --- доля показов не лучшего варианта.

\begingroup\footnotesize\setlength{\tabcolsep}{4pt}
\begin{xltabular}{\linewidth}{@{}lrrr@{}}
\toprule\noalign{}
policy\_name & reward & regret & suboptimal\_share \\
\midrule\noalign{}
\endhead
\bottomrule\noalign{}
\endlastfoot
thompson\_sampling & 213,4 & \textbf{36,6 ± 14,1} & 0,64 \\
epsilon\_greedy & 211,6 & 38,4 ± 30,0 & 0,63 \\
ucb1 & 201,0 & 49,1 ± 12,2 & 0,77 \\
fixed\_ab & 197,5 & 52,6 ± 12,0 & 0,80 \\
\end{xltabular}
\endgroup

Все три MAB-алгоритма опережают \texttt{fixed\_ab} на 14--16 пунктов regret. Лидирует TS; $\varepsilon$-greedy близок, но с большим разбросом по seed. UCB1 отстаёт из-за консервативного exploration при близких CTR и коротком горизонте.

\textbf{Устойчивость результата по seed.} Эксперимент E1 повторялся с 20 независимыми начальными значениями генератора (seed): при каждом seed среда и политики ведут себя по-разному, поэтому regret слегка меняется от прогона к прогону. Чтобы проверить, не случайна ли разница между TS и \texttt{fixed\_ab}, для каждого seed вычислена разность regret(\texttt{fixed\_ab}) $-$ regret(TS) --- на сколько кликов TS «недобрал» меньше, чем фиксированное разбиение. Средняя разность --- 16 кликов за горизонт T = 5000. Доверительный интервал 95\% для этой средней (по формуле для 20 пар наблюдений) --- примерно от 8 до 24 кликов; bootstrap-пересчёт по тем же 20 seed даёт [7,5; 24,5]. Оба интервала не содержат ноль: выигрыш TS воспроизводится при смене seed, а не объясняется одним удачным прогоном. Это проверка устойчивости метрик блока I (онлайн-оптимизация), а не проверка гипотезы о превосходстве B над A; для последней нужны процедуры блока II (§3.6.2).

\textbf{Эксперимент E6 (RQ4)} на той же среде исследует зависимость regret от $\varepsilon$ в $\varepsilon$-greedy при $\varepsilon$ $\in$ \{0,05; 0,10; 0,20\}, по 10 seed. Regret монотонно падает с 55,4 до 36,0, suboptimal\_share --- с 0,83 до 0,63 (табл. 3.3). Слишком малое $\varepsilon$ ведёт себя почти как greedy и проигрывает TS при том же горизонте.

\textbf{Таблица 3.3 --- E6: абляция $\varepsilon$-greedy (среднее, 10 seed, T = 5000)}

\begingroup\footnotesize\setlength{\tabcolsep}{4pt}
\begin{xltabular}{\linewidth}{@{}rrrr@{}}
\toprule\noalign{}
$\varepsilon$ & reward & regret & suboptimal\_share \\
\midrule\noalign{}
\endhead
\bottomrule\noalign{}
\endlastfoot
0,05 & 194,60 & 55,40 & 0,83 \\
0,10 & 204,00 & 46,00 & 0,76 \\
0,20 & 214,00 & \textbf{36,00} & 0,63 \\
\end{xltabular}
\endgroup

\textbf{Дополнение к E1 --- насколько CTR рук отличаются друг от друга.} В базовом E1 CTR пяти рук близки (3--5\%), но не задают крайних случаев. Дополнительно проверены два синтетических сценария с тем же T = 5000 и 20 seed; в коде они названы \texttt{small\_gap} и \texttt{large\_gap}, где \emph{gap} --- разрыв по CTR между лучшей и остальными руками.

\textit{Малый разрыв} (\texttt{small\_gap}): все пять рук почти равны по качеству (CTR от 4,6\% до 5,4\%, шаг около 0,2 п.п.). Алгоритму трудно быстро выделить лидера, поэтому адаптивное назначение мало чем помогает: regret у всех политик близок (21--26), TS почти не отличается от \texttt{fixed\_ab}.

\textit{Большой разрыв} (\texttt{large\_gap}): один вариант явно лучший (CTR 15\%), у остальных четыёх --- 1--4\%. Здесь бандит быстро перенаправляет трафик к лидеру: regret TS --- 28,3 клика, у \texttt{fixed\_ab} --- 505,8, потому что половина показов по-прежнему уходит на слабые руки (табл. 3.4, рис. 3.4).

Для обоих сценариев построены bootstrap-доверительные интервалы разности regret(политика) $-$ regret(\texttt{fixed\_ab}) по 20 seed. При малом разрыве интервалы перекрывают ноль --- статистически убедительного отрыва нет. При большом разрыве интервалы для TS и $\varepsilon$-greedy не включают ноль: преимущество бандита согласуется с графиком.

\textbf{Таблица 3.4 --- Сценарии с малым и большим разрывом CTR (средний regret, T = 5000, 20 seed)}

Строки: бейзлайн E1 (CTR 3--5\%); \texttt{small\_gap} --- руки почти равны; \texttt{large\_gap} --- один явный лидер.

\begingroup\footnotesize\setlength{\tabcolsep}{4pt}
\begin{xltabular}{\linewidth}{@{}lrrrr@{}}
\toprule\noalign{}
scenario & fixed\_ab & $\varepsilon$-greedy & UCB1 & Thompson Sampling \\
\midrule\noalign{}
\endhead
\bottomrule\noalign{}
\endlastfoot
бейзлайн (E1) & 52,6 & 38,4 & 49,1 & \textbf{36,6} \\
малый разрыв (\texttt{small\_gap}) & 25,9 & 26,8 & 24,5 & \textbf{21,0} \\
большой разрыв (\texttt{large\_gap}) & 505,8 & 103,0 & 217,2 & \textbf{28,3} \\
\end{xltabular}
\endgroup

\begin{figure}[htbp]
\centering
\includegraphics[width=0.92\textwidth]{fig_3_4_synthetic_gap_regret.png}
\caption{Рисунок 3.4 — Regret при малом и большом разрыве CTR между руками}
\label{fig:рисунок_3_4___regret_по_сценариям_разбро}
\end{figure}
\begin{figure}[htbp]
\centering
\includegraphics[width=0.92\textwidth]{fig_3_6_epsilon_regret.png}
\caption{Рисунок 3.6 — Абляция $\varepsilon$ (E6)}
\label{fig:рисунок_3_6___абляция_$\varepsilon$__e6_}
\end{figure}

\textbf{Эксперимент E8b (RQ5)} --- контекстный бандит на синтетике: пять рук, четырёхмерный вектор признаков пользователя (контекст), CTR каждой руки --- сигмоида от скалярного произведения. T = 10 000, 20 seed. Сравниваются \texttt{linucb}, \texttt{thompson\_sampling} и \texttt{fixed\_ab}. При информативном контексте LinUCB даёт наименьший regret (26,6 ± 30,0) против 28,3 у TS и 49,2 у \texttt{fixed\_ab} (табл. 3.5). TS чаще выбирает неоптимальную руку (\texttt{suboptimal\_share} 0,86), поскольку не использует признаки.

\textbf{Таблица 3.5 --- E8b: контекстный бандит, синтетика (среднее ± ст. откл., 20 seed, T = 10 000)}

\begingroup\footnotesize\setlength{\tabcolsep}{4pt}
\begin{xltabular}{\linewidth}{@{}lrrr@{}}
\toprule\noalign{}
policy\_name & reward & regret & suboptimal\_share \\
\midrule\noalign{}
\endhead
\bottomrule\noalign{}
\endlastfoot
linucb & 1973,4 & \textbf{26,6 ± 30,0} & 0,78 \\
thompson\_sampling & 1971,7 & 28,3 ± 33,2 & 0,86 \\
fixed\_ab & 1950,8 & 49,2 ± 38,2 & 0,80 \\
\end{xltabular}
\endgroup

\textbf{Стабильность LinUCB vs fixed (E8b).} Средний regret LinUCB ниже fixed на 22,6 пункта, но разброс велик ($\sigma$ $\approx$ 30--38). Bootstrap по 20 seed: 95\%-ДИ разности (fixed - LinUCB) $\approx$ [-5; 50] --- включает 0; превосходство LinUCB над fixed по regret статистически не подтверждено, хотя точечная оценка благоприятна. Сравнение LinUCB с TS (26,6 vs 28,3) на этом объёме seed также не различимо. E8b демонстрирует потенциал контекстного бандита при информативных признаках.

\begin{figure}[htbp]
\centering
\includegraphics[width=0.92\textwidth]{fig_3_1_e8b_regret_curves.png}
\caption{Рисунок 3.1 — E8b: средний cumulative regret (контекстный бандит, синтетика, 20 seed, T = 10 000)}
\label{fig:рисунок_3_1___e8b__средний_cumulative_re}
\end{figure}

В совокупности E1, E6 и E8b показывают: адаптивный онлайн-эксперимент с MAB накапливает больше кликов, чем фиксированное A/B-разбиение; выигрыш усиливается при большом разрыве CTR (табл. 3.4). В E8b наблюдается тенденция в пользу LinUCB при информативном контексте, но bootstrap не подтверждает значимого отрыва от fixed и TS. Это ответ на вопрос блока I --- «приносит ли пользу» --- и не ответ на вопрос статистической значимости для проверки гипотезы о превосходстве B; последний проверяется в §3.6.2.

\textbf{Дополнительные абляции (notebook 03).} На той же синтетике E1 проверена чувствительность Thompson Sampling к априорным параметрам Beta($\alpha$, $\beta$) в сглаживании CTR \texttt{($\alpha$,\ $\beta$)\ $\in$\ \{(1,1),\ (1,5),\ (5,1),\ (10,10)\}} --- разброс mean regret небольшой (рис. 3.7). На OBD в batch-режиме (10k событий) сравнивались \texttt{batch\_size\ $\in$\ \{1,\ 50,\ 200,\ 1000\}}: при укрупнении пакета regret \texttt{fixed\_ab} и $\varepsilon$-greedy сходятся, что иллюстрирует компромисс между частотой обновления и шумом оценки (рис. 3.8).

\begin{figure}[htbp]
\centering
\includegraphics[width=0.92\textwidth]{fig_3_7_ts_priors_regret.png}
\caption{Рисунок 3.7 — Абляция априорных параметров Beta Thompson Sampling}
\label{fig:рисунок_3_7___абляция_priors_thompson_sa}
\end{figure}
\begin{figure}[htbp]
\centering
\includegraphics[width=0.92\textwidth]{fig_3_8_batch_size_regret.png}
\caption{Рисунок 3.8 — Абляция batch_size (OBD, batch-режим)}
\label{fig:рисунок_3_8___абляция_batch_size__obd__b}
\end{figure}

\subsubsection{Блок II: статистическая значимость и интерпретация логов (E4, E5, E11, E14)}\label{ux431ux43bux43eux43a-ii-ux441ux442ux430ux442ux438ux441ux442ux438ux447ux435ux441ux43aux430ux44f-ux437ux43dux430ux447ux438ux43cux43eux441ux442ux44c-ux438-ux438ux43dux442ux435ux440ux43fux440ux435ux442ux430ux446ux438ux44f-ux43bux43eux433ux43eux432-e4-e5-e11-e14}

Блок I показал измеримую пользу адаптивного назначения. Блок II отвечает на продуктовый вопрос, который обычно следует за ним: можно ли по логам того же эксперимента заявить статистически значимый эффект и принять решение о превосходстве варианта B? Литература (§2.2.2) предупреждает: наивные A/B-процедуры здесь часто некорректны. Ниже проверяются три типичных нарушения и известные обходные пути.

Больше кликов во время эксперимента не равно статистически доказанному превосходству B над A. Для проверки гипотезы нужен отдельный статистический вывод, и к логам бандита его нельзя применять наивно.

Продуктовая гипотеза: CTR варианта B выше, чем у A. В экспериментах E4, E11 и E14 используются два сценария. В сценарии \textbf{null} CTR\_A = CTR\_B --- варианты одинаковы, гипотеза «B лучше» ложна; хороший тест не должен отвергать H$_0$. В сценарии \textbf{effect} CTR\_B \textgreater{} CTR\_A --- эффект есть; хороший тест должен отвергать H$_0$ с высокой вероятностью (мощность). Rejection rate при null --- доля прогонов, где тест ошибочно заявил «эффект есть»; при уровне значимости $\alpha$ = 5\% она должна быть около 5\%.

Обычный z-критерий A/B предполагает фиксированное разбиение трафика до начала эксперимента и условную независимость наблюдений. Бандит нарушает оба условия: перекладывает трафик по промежуточным результатам. Отдельная проблема возникает при непрерывном мониторинге фиксированного A/B: если смотреть p-value на промежуточных долях горизонта и останавливать эксперимент при первом «значимом» результате, наивный критерий завышает ошибку I рода даже при честном разбиении 50/50 \cite{ref19}. Для этого случая в литературе предлагают последовательный групповой дизайн с границами O'Brien--Fleming (OBF, см. сокращения) \cite{ref19}; адаптивное назначение трафика бандитом --- другая проблема, и одна процедура не заменяет другую.

\textbf{Эксперимент E4 (RQ3.1).} Monte Carlo: 200 прогонов, T = 20 000, $\alpha$ = 0,05. Логи генерируются \texttt{fixed\_ab} или Thompson Sampling; к обоим применяется один наивный z-критерий. Горизонт 20 000 выбран для осмысленной мощности при effect +1 п.п. CTR (5\% vs 6\%).

\textbf{Таблица 3.6 --- E4: наивный A/B на фиксированном и адаптивном назначении (T = 20 000)}

Столбцы: \texttt{scenario} --- null (нет эффекта) или effect; \texttt{policy} --- политика генерации логов; \texttt{CTR\_A}, \texttt{CTR\_B} --- истинные CTR контрольного и тестового вариантов; \texttt{rejection\_rate} --- доля прогонов с отклонением H$_0$; \texttt{target} --- целевое значение (0,05 при null, 0,80 при effect).

\begingroup\footnotesize\setlength{\tabcolsep}{4pt}
\begin{xltabular}{\linewidth}{@{}llrrrr@{}}
\toprule\noalign{}
scenario & policy & CTR\_A & CTR\_B & rejection\_rate & target \\
\midrule\noalign{}
\endhead
\bottomrule\noalign{}
\endlastfoot
null & fixed\_ab & 0,05 & 0,05 & \textbf{0,030} & 0,05 \\
null & thompson\_sampling & 0,05 & 0,05 & \textbf{0,145} & 0,05 \\
effect & fixed\_ab & 0,05 & 0,06 & \textbf{0,850} & 0,80 \\
effect & thompson\_sampling & 0,05 & 0,06 & 0,550 & 0,80 \\
\end{xltabular}
\endgroup

\textbf{Таблица 3.7 --- E4: мощность при росте эффекта (fixed\_ab, T = 20 000, 200 trials)}

\begingroup\footnotesize\setlength{\tabcolsep}{4pt}
\begin{xltabular}{\linewidth}{@{}rrrr@{}}
\toprule\noalign{}
CTR\_B & ΔCTR & rejection\_rate fixed\_ab & rejection\_rate TS (наивный) \\
\midrule\noalign{}
\endhead
\bottomrule\noalign{}
\endlastfoot
0,06 & +1 п.п. & 0,850 & 0,550 \\
0,08 & +3 п.п. & \textbf{1,000} & 0,975 \\
\end{xltabular}
\endgroup

При слабом эффекте +1 п.п. (табл. 3.7) \texttt{fixed\_ab} обнаруживает его в 85\% прогонов, наивный тест на логах TS --- только в 55\%. При крупном +3 п.п. оба подхода почти всегда срабатывают (100\% и 97,5\%); на слабом эффекте бандит отстаёт из-за \emph{перекоса групп} --- неравномерного распределения показов между A и B после адаптивного назначения (см. E5 ниже).

При отсутствии эффекта (null, табл. 3.6) \texttt{fixed\_ab} с наивным z-критерием даёт rejection rate 3,0\% --- чуть ниже номинальных 5\% из-за конечного числа Monte Carlo-прогонов (200), дискретности кликов и одностороннего теста CTR\_B \textgreater{} CTR\_A. На логах Thompson Sampling тот же наивный тест ошибочно «находит эффект» в 14,5\% прогонов (E4) и в 10\% (E11, 100 trials) --- заметно чаще, чем допускает $\alpha$ = 5\%. При реальном эффекте +1 п.п. мощность \texttt{fixed\_ab} --- 85\%, у TS с наивным тестом --- 55\%.

\textbf{Эксперимент E5} наглядно показывает, откуда берётся перекос. Взят один прогон (T = 5000, seed = 42) при истинных CTR 4\% и 6\% --- эффект есть, но группы после адаптивного назначения сильно неравны (табл. 3.8).

\textbf{Таблица 3.8 --- E5: один прогон, неравные группы A и B (T = 5000, seed = 42)}

Столбцы: \texttt{n\_A}, \texttt{n\_B} --- число показов контрольного и тестового вариантов; при честном 50/50 ожидалось бы около 2500 в каждой группе.

\begingroup\footnotesize\setlength{\tabcolsep}{4pt}
\begin{xltabular}{\linewidth}{@{}lrrrl@{}}
\toprule\noalign{}
policy & n\_A & n\_B & p\_value & вывод \\
\midrule\noalign{}
\endhead
\bottomrule\noalign{}
\endlastfoot
fixed\_ab & 2443 & 2557 & 0,0018 & эффект \\
thompson\_sampling & 443 & 4557 & 0,0235 & эффект \\
epsilon\_greedy & 4699 & 301 & 0,4595 & нет эффекта \\
ucb1 & 2221 & 2779 & 0,1607 & нет эффекта \\
\end{xltabular}
\endgroup

У Thompson Sampling 90\% показов (4557 из 5000) ушло варианту B --- z-критерий «видит» высокий CTR у B, хотя это следствие перераспределения трафика, а не корректного A/B-дизайна. У $\varepsilon$-greedy перекос обратный (4699 vs 301): тест не замечает реальный эффект. Такие искажения разбирает Xiang et al.~\cite{ref21}; E4 измеряет их последствия в rejection rate и мощности, E11 проверяет исправление через IPS, E14 --- отдельную проблему подглядывания при фиксированном A/B.

Практический вывод из E4--E5: бандит полезен для накопления кликов (блок I), но не заменяет эксперимент, заточенный под проверку гипотезы. Если нужно обоснованно заявить «B лучше A», надёжнее заранее запланировать A/B с \texttt{fixed\_ab} --- в нашей постановке при T = 20 000 и эффекте +1 п.п. ошибка I рода $\approx$ 3\%, мощность $\approx$ 85\%. Если отдельный A/B уже невозможен и остались только bandit-логи, применяют IPS-взвешенный статистический вывод; его проверяет E11.

\textbf{Эксперимент E11 (RQ3.2)} моделирует второй случай: продукт уже собирал логи под Thompson Sampling, выделить новый трафик под чистый A/B нельзя, но нужно постфактум проверить CTR\_B \textgreater{} CTR\_A. На синтетике генерируются bandit-логи с propensity; сравниваются три процедуры: наивный z-критерий на TS-логах (как в E4, ожидаемо с завышенной ошибкой I рода), IPS-взвешенный вывод (коррекция с учётом propensity) и z-критерий на эталонных логах \texttt{fixed\_ab} (что было бы при идеальном отдельном A/B). Monte Carlo: 100 прогонов, T = 20 000, $\alpha$ = 0,05; сценарии null (5\% / 5\%) и effect (+1, +3, +5 п.п. к CTR\_B). Для стандартной ошибки IPS в \texttt{run\_all} включён bootstrap по строкам лога (200 ресэмплов).

\textbf{Таблица 3.9 --- E11: ошибка I рода и слабый эффект (+1 п.п., T = 20 000)}

Столбцы: \texttt{method} --- процедура вывода; \texttt{rejection\_rate} --- доля отклонений H$_0$; \texttt{target} --- целевое значение.

\begingroup\footnotesize\setlength{\tabcolsep}{4pt}
\begin{xltabular}{\linewidth}{@{}llrr@{}}
\toprule\noalign{}
scenario & method & rejection\_rate & target \\
\midrule\noalign{}
\endhead
\bottomrule\noalign{}
\endlastfoot
null & naive\_ab\_full\_ts & 0,100 & 0,05 \\
null & naive\_ab\_fixed\_only & 0,030 & 0,05 \\
null & ips\_weighted\_ab & 0,050 & 0,05 \\
effect (+1 п.п.) & naive\_ab\_full\_ts & 0,510 & 0,80 \\
effect (+1 п.п.) & naive\_ab\_fixed\_only & 0,850 & 0,80 \\
effect (+1 п.п.) & ips\_weighted\_ab & 0,240 & 0,80 \\
\end{xltabular}
\endgroup

\textbf{Таблица 3.10 --- E11: мощность IPS при росте эффекта (\texttt{ips\_weighted\_ab}, T = 20 000)}

\begingroup\footnotesize\setlength{\tabcolsep}{4pt}
\begin{xltabular}{\linewidth}{@{}rlrr@{}}
\toprule\noalign{}
ΔCTR & CTR\_A → CTR\_B & rejection\_rate IPS & rejection\_rate fixed\_ab \\
\midrule\noalign{}
\endhead
\bottomrule\noalign{}
\endlastfoot
+1 п.п. & 5\% → 6\% & 0,240 & 0,850 \\
+3 п.п. & 5\% → 8\% & 0,780 & 1,000 \\
+5 п.п. & 5\% → 10\% & 0,970 & 1,000 \\
\end{xltabular}
\endgroup

IPS восстанавливает контроль ошибки I рода на адаптивных логах (5\% против 10\% у наивного теста). На слабом +1 п.п. мощность IPS низкая (24\%): адаптивный лог даёт мало информации о редко показываемой руке (низкий ESS). При +3 п.п. IPS находит эффект в 78\% прогонов, при +5 п.п. --- в 97\%. Для проверки слабой гипотезы нужен отдельный A/B; IPS --- постфактумный инструмент при заметном эффекте.

\textbf{Эксперимент E14 (RQ3.3)} отделяет подглядывание за промежуточными результатами при фиксированном разбиении трафика от адаптивного назначения бандита. Пять процедур: \texttt{fixed\_horizon}, \texttt{naive\_peek} (4 промежуточных среза), последовательный групповой дизайн \textbf{OBF} \cite{ref33}, корректный мониторинг \textbf{mSPRT} \cite{ref19} ($\tau$ = 0,02). Monte Carlo: 200 trials, T = 20 000, $\alpha$ = 0,05.

\textbf{Таблица 3.11 --- E14: подглядывание и последовательный статистический вывод (T = 20 000, 200 trials)}

Столбцы: \texttt{method} --- процедура вывода; \texttt{rejection\_rate} --- доля отклонений H$_0$; \texttt{mean\_stop\_fraction} --- средняя доля горизонта, на которой эксперимент остановлен. OBF --- 4 looks (25/50/75/100\%); mSPRT --- непрерывный мониторинг ($\tau$ = 0,02, burn-in 100 пар).

\begingroup\footnotesize\setlength{\tabcolsep}{4pt}
\begin{xltabular}{\linewidth}{@{}>{\raggedright\arraybackslash}p{0.10\linewidth}>{\raggedright\arraybackslash}p{0.14\linewidth}>{\raggedright\arraybackslash}p{0.34\linewidth}>{\raggedleft\arraybackslash}r>{\raggedleft\arraybackslash}r>{\raggedleft\arraybackslash}r@{}}
\toprule\noalign{}
scenario & policy & method & rejection\_rate & mean\_stop\_fraction & target \\\\
\midrule\noalign{}
\endhead
\bottomrule\noalign{}
\endlastfoot
null & fixed\_ab & fixed\_horizon & 0,030 & 1,00 & 0,05 \\
null & fixed\_ab & naive\_peek & \textbf{0,110} & 0,95 & 0,05 \\
null & fixed\_ab & group\_sequential\_obf & \textbf{0,040} & 1,00 & 0,05 \\
null & fixed\_ab & always\_valid\_msprt & \textbf{0,020} & 0,98 & 0,05 \\
null & thompson\_sampling & fixed\_horizon & 0,145 & 1,00 & 0,05 \\
null & thompson\_sampling & naive\_peek & \textbf{0,210} & 0,89 & 0,05 \\
null & thompson\_sampling & group\_sequential\_obf & \textbf{0,145} & 0,99 & 0,05 \\
null & thompson\_sampling & always\_valid\_msprt & 0,005 & 1,00 & 0,05 \\
effect & fixed\_ab & fixed\_horizon & 0,850 & 1,00 & 0,80 \\
effect & fixed\_ab & naive\_peek & 0,900 & 0,61 & 0,80 \\
effect & fixed\_ab & group\_sequential\_obf & 0,855 & 0,80 & 0,80 \\
effect & fixed\_ab & always\_valid\_msprt & 0,585 & 0,74 & 0,80 \\
effect & thompson\_sampling & fixed\_horizon & 0,550 & 1,00 & 0,80 \\
effect & thompson\_sampling & naive\_peek & 0,660 & 0,69 & 0,80 \\
effect & thompson\_sampling & group\_sequential\_obf & 0,560 & 0,93 & 0,80 \\
effect & thompson\_sampling & always\_valid\_msprt & 0,010 & 1,00 & 0,80 \\
\end{xltabular}
\endgroup

На логах \texttt{fixed\_ab} наивное подглядывание при null поднимает rejection rate с 3\% до 11\%: многократные просмотры промежуточных p-value без поправки раздувают число ложных срабатываний. Обе корректные последовательные процедуры восстанавливают контроль: OBF --- 4\%, mSPRT --- 2\%. При effect OBF сохраняет мощность, близкую к фиксированному горизонту (85,5\% против 85\%); mSPRT останавливает эксперимент раньше (mean\_stop\_fraction $\approx$ 0,74), но при слабом +1 п.п. детектирует эффект только в 58,5\% прогонов --- ниже целевых 80\%. Это согласуется с Johari et al.~\cite{ref19}: для Bernoulli при умеренном $\alpha$ формула mSPRT приближённая, а мощность зависит от заранее выбранного $\tau$; OBF с четырьмя промежуточными срезами проще в реализации и в данной постановке мощнее.

На логах Thompson Sampling ни OBF, ни mSPRT не заменяют IPS и не исправляют адаптивное назначение. OBF при null даёт те же 14,5\%, что наивный тест в конце горизонта: разбиение не фиксировано, и последовательный групповой дизайн не рассчитан на назначение, зависящее от данных. mSPRT ведёт себя иначе --- rejection rate при null падает до 0,5\%, мощность при effect --- до 1\% --- не потому что вывод стал корректным, а потому что при перекосе групп эффективное число пар мало и тест теряет чувствительность. Johari et al.~\cite{ref19} оговаривают, что корректный мониторинг при непрерывном наблюдении (always-valid, см. сокращения) распространяется на фиксированное назначение; bandit-логи --- отдельная задача (E4 → отдельный A/B; E11 → IPS).

\textbf{Итог блока II.} Три независимые проблемы интерпретации логов онлайн-эксперимента покрыты эмпирически: (1) адаптивное назначение --- E4, E5 → отдельный fixed A/B или IPS (E11); (2) постфактумный анализ bandit-логов --- E11; (3) подглядывание при фиксированном разбиении --- E14 → OBF / mSPRT. Ни одна последовательная процедура не заменяет IPS на логах бандита; cumulative CTR из блока I не заменяет p-value из блока II.

\textbf{Дополнение E15 --- CUPED при фиксированном разбиении (синтетика).} Отдельный Monte Carlo на fixed 50/50 A/B с коррелированной доэкспериментной ковариатой (\texttt{pre\_reward}): 100 trials, T = 20 000, $\alpha$ = 0,05. Сравниваются наивный z-тест и CUPED \cite{ref7} на скорректированных исходах (табл. 3.1; \texttt{outputs/cuped\_power/e15\_summary.csv}). При null оба метода дают rejection rate 5\% (100 trials, номинальный $\alpha$ = 5\%). При effect +1 п.п. CUPED даёт мощность 85\% против 84\% у наивного A/B; средний SE у CUPED ниже (\textasciitilde0,00316 против \textasciitilde0,00320). При +3 п.п. оба достигают 100\% мощности. На синтетике CUPED даёт скромное снижение дисперсии; на OBD ковариаты ограничены (§4.3).

\subsubsection{Интерпретация логов OBD: офлайн-оценка (E2, E2b/c, E9)}\label{ux438ux43dux442ux435ux440ux43fux440ux435ux442ux430ux446ux438ux44f-ux43bux43eux433ux43eux432-obd-ux43eux444ux43bux430ux439ux43d-ux43eux446ux435ux43dux43aux430-e2-e2bc-e9}

На синтетике блоки I и II проверялись \emph{отдельными} экспериментами: в блоке I измеряли клики и regret (E1, E6, E8b), в блоке II --- корректность статистического вывода (E4, E5, E11, E14), не смешивая эти метрики в одном заключении. На реальных логах OBD вопрос другой: эксперимент уже завершён под конкретной политикой логирования, новый трафик не выделяют, а сравнивают кандидатов постфактум --- через IPS/SNIPS replay-OPE \cite{ref24}, \cite{ref26}. OPE не заменяет блок II: по SNIPS нельзя получить p-value о превосходстве B над A. Главные диагностические метрики --- доля совпадений с логом, ESS и стабильность SNIPS между seed; сравнение адаптивных политик по максимальному SNIPS при низкой доле совпадений не интерпретируется как ранжирование политик. Дополнительно в пайплайне доступен bootstrap-ДИ для SNIPS (ресэмплирование строк лога с возвращением, \texttt{-\/-bootstrap} в \texttt{compare\_ab\_vs\_bandits} / \texttt{obd\_streaming\_ope}); в \texttt{run\_all} включён для E2/E9/E13.

\textbf{Эксперимент E2 (RQ2)} --- IPS/SNIPS replay-OPE на \texttt{random/all}: 10 000 событий, 80 рук, логирование при равномерном назначении. Четыре политики-кандидата, 20 seed. Доля совпадений с логом $\approx$ 1,2\% у всех политик (ожидаемо при 80 руках); ESS $\approx$ 120--124 (табл. 3.12). SNIPS-оценки CTR лежат на уровне 10⁻³ с высокой дисперсией между seed; 38 кликов на 10 000 показов недостаточно для уверенного ранжирования. OPE демонстрирует пайплайн, но не выбирает «лучший бандит для продакшена».

\textbf{Таблица 3.12 --- E2: IPS/SNIPS replay-OPE (среднее, 20 seed)}

Столбцы: \texttt{acceptance\_rate} --- доля совпадений с логом (рука политики-кандидата совпадает с залогированной); \texttt{snips\_estimate} --- оценка CTR; \texttt{ESS} --- эффективный размер выборки.

\begingroup\footnotesize\setlength{\tabcolsep}{4pt}
\begin{xltabular}{\linewidth}{@{}lrrr@{}}
\toprule\noalign{}
policy\_name & acceptance\_rate & snips\_estimate & ESS \\
\midrule\noalign{}
\endhead
\bottomrule\noalign{}
\endlastfoot
thompson\_sampling & 0,0121 & 0,0049 & 120,6 \\
fixed\_ab & 0,0123 & 0,0029 & 123,4 \\
epsilon\_greedy & 0,0124 & 0,0004 & 123,6 \\
ucb1 & 0,0122 & 0,0000 & 122,0 \\
\end{xltabular}
\endgroup

\textbf{Эксперимент E2b/c (RQ6)} расширяет E2 на четыре подвыборки: \texttt{random/all}, \texttt{bts/all}, \texttt{random/men}, \texttt{random/women}, 20 seed. На \texttt{bts/all} ESS для \texttt{fixed\_ab} падает до 58 против 123 на \texttt{random/all} --- весовая оценка деградирует при адаптивном логировании (табл. 3.13, рис. 3.2). Максимальный mean SNIPS по подвыборкам приведён как иллюстрация дисперсии оценщика, не как рекомендация политики.

\begin{figure}[htbp]
\centering
\includegraphics[width=0.92\textwidth]{fig_3_2_ope_ess_matrix.png}
\caption{Рисунок 3.2 — E2b/c: ESS по политике и подвыборке OBD}
\label{fig:рисунок_3_2___e2b_c__ess_по_политике_и_п}
\end{figure}

\textbf{Таблица 3.13 --- E2b/c: максимальный mean SNIPS (иллюстрация; не ранжирование политик)}

\begingroup\footnotesize\setlength{\tabcolsep}{4pt}
\begin{xltabular}{\linewidth}{@{}lrlr@{}}
\toprule\noalign{}
dataset & max mean SNIPS & policy\_name & ESS \\
\midrule\noalign{}
\endhead
\bottomrule\noalign{}
\endlastfoot
random\_all & 0,0049 & thompson\_sampling & 120,6 \\
bts\_all & 0,0024 & fixed\_ab & 57,6 \\
random\_men & 0,0147 & ucb1 & 272,0 \\
random\_women & 0,0041 & thompson\_sampling & 217,8 \\
\end{xltabular}
\endgroup

Разведочные batch-прогоны E3 и E8 на OBD (не IPS/SNIPS replay-OPE, без контрфактической оценки для LinUCB) вынесены в приложение А --- в основной текст не включены.

\textbf{Эксперимент E9} масштабирует OPE: E9a --- полный \texttt{random/all} (1 374 327 событий, 4 768 кликов); E9b --- первый миллион \texttt{bts/all}, 10 seed, потоковая обработка. ESS вырос примерно в 140 раз (123 → 17 152 для \texttt{fixed\_ab}), доля совпадений с логом остаётся $\approx$ 1,25\%. SNIPS между политиками различаются на уровне 10⁻³ при заметной дисперсии между seed (табл. 3.14--3.15); устойчивого ранжирования не получено.

\textbf{Таблица 3.14 --- E9a: полный \texttt{random/all} (среднее ± ст. откл., 10 seed)}

\begingroup\footnotesize\setlength{\tabcolsep}{4pt}
\begin{xltabular}{\linewidth}{@{}lrrr@{}}
\toprule\noalign{}
policy\_name & snips\_estimate & ESS & acceptance\_rate \\
\midrule\noalign{}
\endhead
\bottomrule\noalign{}
\endlastfoot
fixed\_ab & 0,00360 ± 0,00022 & 17 152 & 1,25\% \\
thompson\_sampling & 0,00343 ± 0,00023 & 17 217 & 1,25\% \\
epsilon\_greedy & 0,00278 ± 0,00021 & 17 603 & 1,28\% \\
ucb1 & 0,00272 ± 0,00000 & 17 657 & 1,28\% \\
\end{xltabular}
\endgroup

\textbf{Таблица 3.15 --- E9b: bts/all (среднее ± ст. откл., 10 seed)}

\begingroup\footnotesize\setlength{\tabcolsep}{4pt}
\begin{xltabular}{\linewidth}{@{}lrrr@{}}
\toprule\noalign{}
policy\_name & snips\_estimate & ESS & acceptance\_rate \\
\midrule\noalign{}
\endhead
\bottomrule\noalign{}
\endlastfoot
thompson\_sampling & 0,0040 ± 0,0008 & 6 006 & 1,25\% \\
fixed\_ab & 0,0039 ± 0,0007 & 6 023 & 1,25\% \\
epsilon\_greedy & 0,0023 ± 0,0003 & 2 888 & 0,36\% \\
ucb1 & 0,0016 ± 0,0000 & 2 596 & 0,26\% \\
\end{xltabular}
\endgroup

\textbf{Таблица 3.16 --- сравнение масштаба OPE (\texttt{fixed\_ab})}

\begingroup\footnotesize\setlength{\tabcolsep}{4pt}
\begin{xltabular}{\linewidth}{@{}>{\raggedright\arraybackslash}p{0.22\linewidth}>{\raggedleft\arraybackslash}r>{\raggedleft\arraybackslash}r>{\raggedleft\arraybackslash}r>{\raggedright\arraybackslash}p{0.33\linewidth}@{}}
\toprule\noalign{}
Этап & событий & кликов & ESS & примечание \\\\
\midrule\noalign{}
\endhead
\bottomrule\noalign{}
\endlastfoot
E2, random/all & 10k & 38 & 123 & доля совпадений $\approx$ 1,2\% \\
E9a, random/all & 1,37M & 4 768 & 17 152 & SNIPS \textasciitilde0,0036, высокая дисперсия \\
E2b/c, bts/all & 10k & 42 & 58 & адаптивное логирование \\
E9b, bts/all & 1M & 3 500 & 6 023 & SNIPS \textasciitilde0,004, доля совпадений $\approx$ 1,25\% у fixed\_ab \\
\end{xltabular}
\endgroup

Цепочка E2 → E2b/c → E9 показывает: replay-OPE-пайплайн работает; ранжирование на малых выборках хрупко; масштаб улучшает ESS, но не даёт устойчивого ранжирования на low-CTR логах с 80 руками; политика логирования критична.

\subsubsection{Классическая A/B-постановка (E12, E13, RQ7)}\label{ux43aux43bux430ux441ux441ux438ux447ux435ux441ux43aux430ux44f-ab-ux43fux43eux441ux442ux430ux43dux43eux432ux43aux430-e12-e13-rq7}

Стандартный A/B-дизайн с гипотезой CTR\_B \textgreater{} CTR\_A. E12 на синтетике проверяет блок I (клики) и блок II (статистический вывод). E13 --- проекция 80-ручного OBD на две \texttt{item\_id} (272 события, \textasciitilde1 клик на руку); это не классический A/B и не источник выводов о превосходстве варианта. Протокол выбора пары --- приложение А; сравнение постановок --- табл. 3.19.

\textbf{Эксперимент E12.} CTR\_A = 5\%, CTR\_B = 8\%. T = 10 000, 20 seed. \texttt{fixed\_ab} 50/50, TS, $\varepsilon$-greedy ($\varepsilon$ = 0,1).

\textbf{Таблица 3.17 --- E12: онлайн-оптимизация (среднее, 20 seed, T = 10 000)}

Столбцы: \texttt{cumulative\_reward} --- суммарные клики; \texttt{cumulative\_regret} --- недобранные клики; \texttt{suboptimal\_share} --- доля показов не лучшего варианта.

\begingroup\footnotesize\setlength{\tabcolsep}{4pt}
\begin{xltabular}{\linewidth}{@{}>{\raggedright\arraybackslash}p{0.34\linewidth}>{\raggedleft\arraybackslash}r>{\raggedleft\arraybackslash}r>{\raggedleft\arraybackslash}r@{}}
\toprule\noalign{}
policy\_name & cumulative\_reward & cumulative\_regret & suboptimal\_share \\\\
\midrule\noalign{}
\endhead
\bottomrule\noalign{}
\endlastfoot
thompson\_sampling & \textbf{786,9} & \textbf{13,1} & \textbf{0,048} \\
epsilon\_greedy & 772,5 & 27,5 & 0,088 \\
fixed\_ab & 648,4 & 151,7 & 0,501 \\
\end{xltabular}
\endgroup

TS набирает 786,9 кликов при regret 13,1; \texttt{fixed\_ab} --- 648,4 кликов, regret 151,7. Разрыв около 140 кликов на 10 000 показов. Картина совпадает с E1: бандит выигрывает по кликам в попарной постановке.

Второй блок E12 --- статистический вывод на тех же CTR: 100 Monte Carlo trials, $\alpha$ = 0,05. Симулируются логи TS с propensity и эталон \texttt{fixed\_ab}; применяются наивный z-критерий на TS, IPS-взвешенный статистический вывод и наивный z-критерий на fixed-логах.

\textbf{Таблица 3.18 --- E12: статистический вывод при эффекте +3 п.п. (T = 10 000, 100 trials)}

Столбцы: \texttt{method} --- процедура вывода; \texttt{rejection\_rate} --- доля прогонов с отклонением H$_0$; последний столбец --- краткое описание метода.

\begingroup\footnotesize\setlength{\tabcolsep}{4pt}
\begin{xltabular}{\linewidth}{@{}lrl@{}}
\toprule\noalign{}
method & rejection\_rate & смысл \\
\midrule\noalign{}
\endhead
\bottomrule\noalign{}
\endlastfoot
naive\_ab\_fixed\_only & 1,00 & эталон: чистый A/B \\
ips\_weighted\_ab & 0,56 & IPS на логах бандита \\
naive\_ab\_full\_ts & 0,90 & наивный тест на TS-логе \\
\end{xltabular}
\endgroup

Эталонный A/B находит эффект в 100\% прогонов. IPS --- в 56\%. Наивный тест на TS-логе --- в 90\%. Для проверки гипотезы нужен отдельный A/B-дизайн.

\textbf{E13 (кратко).} Проекция OBD на пару items 41/50: в 2-ручной постановке доля совпадений с логом $\approx$ 49\% (против 1,2\% в 80-ручной), ESS $\approx$ 134 при 272 событиях --- иллюстрация интерпретируемости метрик, без выводов о ранжировании.

\textbf{Таблица 3.19 --- многорукий OBD и проекция на пару вариантов: разные вопросы}

\begingroup\footnotesize\setlength{\tabcolsep}{4pt}
\begin{xltabular}{\linewidth}{@{}>{\raggedright\arraybackslash}p{0.22\linewidth}>{\raggedright\arraybackslash}p{0.39\linewidth}>{\raggedright\arraybackslash}p{0.39\linewidth}@{}}
\toprule\noalign{}
 & Multi-arm OBD (E2, E9) & Проекция на пару (E13) \\\\
\midrule\noalign{}
\endhead
\bottomrule\noalign{}
\endlastfoot
Вопрос & Какой алгоритм лучше на 80 товарах? & Метрики на двух вариантах A и B? \\
Руки & 80 & 2 \\
События & 10k -- 1,37M & 272 \\
ESS (\texttt{fixed\_ab}) & 123 -- 17 152 & 134 \\
Связь с A/B & слабая & проекция, не дизайн \\
\end{xltabular}
\endgroup

E12 подтверждает перенос блоков I и II на попарную синтетику. E13 --- методологический мост к продуктовой постановке.

\subsubsection{Вычислительная применимость (P1)}\label{ux432ux44bux447ux438ux441ux43bux438ux442ux435ux43bux44cux43dux430ux44f-ux43fux440ux438ux43cux435ux43dux438ux43cux43eux441ux442ux44c-p1}

\textbf{Эксперимент P1} измеряет фактическое время одного решения (\texttt{select\_arm} + \texttt{update}) при 80 руках, context\_dim = 7, 20 000 вызовов. LinUCB требует около 23 ms на decision, TS --- около 0,05 ms (табл. 3.20). LinUCB не запускался в E9 из-за стоимости. Выбор контекстного бандита --- вопрос не только regret, но и инженерного бюджета.

\textbf{Таблица 3.20 --- P1: время задержки (мкс на decision, 80 рук, 20 000 вызовов)}

Столбец \texttt{total\_decision\_us} --- среднее фактическое время одного вызова \texttt{select\_arm} + \texttt{update} в микросекундах.

\begingroup\footnotesize\setlength{\tabcolsep}{4pt}
\begin{xltabular}{\linewidth}{@{}lr@{}}
\toprule\noalign{}
policy\_name & total\_decision\_us \\
\midrule\noalign{}
\endhead
\bottomrule\noalign{}
\endlastfoot
ucb1 & 2,8 \\
fixed\_ab & 41,6 \\
thompson\_sampling & 51,2 \\
linucb & 23 080 \\
\end{xltabular}
\endgroup

\begin{figure}[htbp]
\centering
\includegraphics[width=0.92\textwidth]{fig_3_3_p1_latency_80arms.png}
\caption{Рисунок 3.3 — P1: время задержки одного решения при 80 руках (log scale)}
\label{fig:рисунок_3_3___p1__время_задержки_одного_}
\end{figure}

\subsection{Краткие выводы по главе 3}\label{ux43aux440ux430ux442ux43aux438ux435-ux432ux44bux432ux43eux434ux44b-ux43fux43e-ux433ux43bux430ux432ux435-3}

Блок I: MAB снижает regret и увеличивает клики (E1, E12); блок II: наивный статистический вывод на bandit-логах некорректен (E4, E11); подглядывание требует последовательных процедур (E14); OPE на OBD --- пайплайн без устойчивого ранжирования (E2, E9). Итоговая продуктовая рекомендация, вклад, ограничения и возможное развитие работы --- в главе 4 (§4.1--4.4).

\subsection{Аудит утечек данных и чувствительность}\label{ux430ux443ux434ux438ux442-ux443ux442ux435ux447ux435ux43a-ux434ux430ux43dux43dux44bux445-ux438-ux447ux443ux432ux441ux442ux432ux438ux442ux435ux43bux44cux43dux43eux441ux442ux44c}

\textbf{Propensity и replay-OPE.} Propensity берётся из поля лога (\texttt{propensity\_score} OBD), а не из политики-кандидата --- утечки «будущей» информации в веса IPS нет \cite{ref26}. Политика-кандидат заморожена при оценке (\texttt{freeze\_policy=True} в \texttt{src/ope/replay.py}): \texttt{policy.update()} не вызывается, оценивается фиксированная π. События перетасовываются per seed (\texttt{shuffle\_events=True}), чтобы снять зависимость от порядка файла. Propensity клипируется снизу (\texttt{propensity\_floor=0.01}, как в \texttt{weighted\_inference.py}). Знаменатель IPS --- число валидных строк лога с положительной propensity, не общий объём файла. Candidate используется только для сравнения действия с логом (доля совпадений с логом). Сравнение SNIPS адаптивных политик при доле совпадений $\approx$ 1/80 не интерпретируется как ранжирование политик.

\textbf{Синтетика.} Шаги среды независимы по seed; CTR фиксированы и не подглядываются при выборе руки. Regret считается относительно оракула с известными CTR --- отдельно от прогонов статистического вывода E4/E11.

\textbf{Статистический вывод.} E4/E11/E14 генерируют свежие логи под заявленную политику; p-value не использует одни и те же данные для назначения и теста вне заявленной модели. IPS-взвешенный статистический вывод использует propensity логирующей политики на том же потоке --- корректная постфактумная постановка \cite{ref30}.

\textbf{batch-режим (E3, E8).} Помечен как невалидный для статистического вывода и IPS/SNIPS replay-OPE (приложение А); в основные выводы не входит.

\textbf{Чувствительность гиперпараметров.} Абляции $\varepsilon$ (E6, рис. 3.6), априорных параметров Beta у TS и \texttt{batch\_size} на OBD (рис. 3.7--3.8) показывают умеренную чувствительность regret к $\varepsilon$ и начальному сглаживанию; вывод блока I по лидерству TS устойчив при $\varepsilon$ $\in$ [0,1; 0,2].

\textbf{Нестационарность.} Дрейф CTR во времени не моделировался --- типичный риск в RecSys \cite{ref16,ref38}; возможное развитие --- §4.3, направление 3.

\subsection{Ограничения}\label{ux43eux433ux440ux430ux43dux438ux447ux435ux43dux438ux44f}

Работа ограничена одним доменом (fashion e-commerce, OBD \cite{ref26}) и одной бизнес-метрикой (CTR). Выводы о применимости MAB и A/B в других продуктах (конверсия, выручка, удержание) требуют отдельной проверки.

\textbf{Данные и масштаб.} Учебная подвыборка OBD (10k событий, \textasciitilde40 кликов) не обеспечивает мощности для ранжирования 80 алгоритмов; даже полный \texttt{random/all} (1,37M событий, \textasciitilde4,8k кликов) не стабилизировал устойчивое ранжирование по SNIPS. Проекция E13 на пару вариантов (272 события) интерпретируема, но не заменяет отдельный A/B.

\textbf{Методология OPE.} Использован SNIPS replay с замороженной политикой-кандидатом, не DR \cite{ref27}; batch E3/E8 --- приложение А. Признаки пользователя (контекст) в OBD строятся MD5-хэшированием в \texttt{scripts/prepare\_open\_bandit.py} --- это не полноценная контекстная OPE для LinUCB на \texttt{random}-логах. Разведочный replay-OPE LinUCB на \texttt{random/women} (10k событий, 20 seed) дал SNIPS $\approx$ 0 при ESS $\approx$ 225 и не включён в табл. 3.12--3.13: при хэш-контексте и низкой доле совпадений с логом оценка не даёт содержательного сравнения с неконтекстными политиками. Выводы о LinUCB опираются на E8b (контролируемая синтетика), а не на OPE по OBD. LinUCB не оценивался в потоковом E9 на полном OBD (P1).

\textbf{Статистический вывод.} mSPRT приближённый \cite{ref19}. IPS: ошибка I рода $\approx$ 5\% (E11, 100 trials, bootstrap SE); наивный TS $\approx$ 10--14,5\% (E11/E4). E11 и OPE используют bootstrap по строкам лога (\texttt{run\_all}); adaptively-weighted CI \cite{ref41} --- направление для диплома. CUPED \cite{ref7} реализован как E15 на синтетике; DR \cite{ref27} --- в §4.3. Весь статистический вывод --- Monte Carlo на синтетике.

\textbf{Нестационарность и GMV.} Дрейф CTR и метрика GMV не измерялись; в работе --- только клики и условная иллюстрация стоимости в §4.2.

\textbf{Реализация.} Время задержки измерено in-process Python без сети. E8b не гарантирует перенос на реальные признаки.

\textbf{Обобщение.} Бандит --- для онлайн-оптимизации (блок I); fixed A/B, IPS или последовательные процедуры --- для проверки гипотезы о превосходстве B (блок II). Возможный продуктовый пилот --- §4.3--4.4.
